\documentclass[11pt]{article}

\usepackage[a4paper,margin=2.6cm]{geometry}
\usepackage{amsmath,amssymb,amsthm}
\usepackage{booktabs}
\usepackage[round,authoryear]{natbib}
\usepackage{microtype}
\usepackage[hidelinks]{hyperref}
\usepackage{listings}
\usepackage{xcolor}

\lstset{
  basicstyle=\ttfamily\footnotesize,
  breaklines=true,
  frame=single,
  framesep=4pt,
  rulecolor=\color{black!25},
  backgroundcolor=\color{black!3},
  showstringspaces=false,
}

\newcommand{\att}{\mathrm{ATT}}
\newcommand{\catt}{\mathrm{CATT}}
\newcommand{\gatt}{\mathrm{GATT}}
\newcommand{\Xb}{\mathbf{X}}

\title{Prognostic--treatment identification in DiD-BCF,\\
       and two repairs to the revision simulation design}
\author{Technical note accompanying the revision}
\date{23 July 2026}

\begin{document}
\maketitle

\begin{abstract}
This note does three things. First, it records two repairs made to the revision
simulation suite: the linearity-degree axis, which previously changed nothing
any estimator could observe, now changes the covariate parameterisation handed
to the estimators (degree~2) and the shape of the treatment-effect surface
(degree~3); and the pre-trend recentring, which annihilated the plain DiD-BCF
estimate, is disabled. Second, and more importantly, it addresses a question
raised by the growing-with-$N$ bias of the posterior correction: whether the
prognostic and treatment forests of DiD-BCF are separately identified. They are
not, and the failure is not an artefact of the revision code. It is present in
the model equation of the submitted paper and in every notebook of the original
simulation study. It does not, however, bite there: refitting the published
specification on the original data-generating process recovers the true effect
to within one per cent at every panel size tested, so the published results
stand. What activates the latent failure is the canonical design the reviewers
asked for, in which an unobserved unit effect correlated with assignment forces
the prognostic forest to split on the treated-group indicator. Third, the note
argues that the failure has a clean interpretation and a clean fix: the
parallel-trends assumption is exactly the restriction that the untreated
outcome has no group-by-time interaction, the current specification does not
encode that restriction, and imposing it additively restores identification.
\end{abstract}

% =============================================================================
\section{The identification problem}
% =============================================================================

\subsection{Statement}

DiD-BCF fits, in the reparameterised form the paper adopts,
\begin{equation}\label{eq:model}
  Y_{it} = \mu(D_i, t, \Xb_{it}) + \tau(\Xb_{it}, k_{it})\, D_{it}
           + \varepsilon_{it},
\end{equation}
where $D_i = \mathbb{I}(G_i \neq \infty)$ is time-invariant ever-treated status,
$D_{it} = D_i \cdot \mathbb{I}(k_{it} \geq 0)$ is the post-treatment indicator,
and $\mu$ and $\tau$ are given independent BART priors. The reparameterisation
is the paper's own contribution and it is a good one: because $\tau$ is
multiplied by $D_{it}$, the estimation problem for $\tau$ is confined to the
post-treatment region and $\tau$ is not asked to learn a structural zero.

The difficulty is on the other side. Write the treatment effect term as a
function of the arguments $\mu$ already has. Since
\[
  D_{it} = \mathbb{I}(D_i = 1)\cdot\mathbb{I}(t \geq g),
\]
the product $\tau(\Xb, k)\,D_{it}$ is a function of $(D_i, t, \Xb)$ and nothing
else. It therefore lies in the domain of $\mu$. For any $c \in [0,1]$ the pair
\begin{equation}\label{eq:flat}
  \bigl(\mu + c\,\tau(\Xb,k)D_{it},\;\; (1-c)\,\tau\bigr)
\end{equation}
produces exactly the same fitted values, the same residuals and the same
likelihood as $(\mu, \tau)$. The data contain no information whatever about
$c$. What separates the two forests is only the prior: in the default
configuration $\mu$ has 200 trees and a diffuse leaf scale while $\tau$ has 50
trees and an aggressive shrinkage prior, so the posterior prefers to place the
effect in $\mu$.

In tree terms the reconstruction is trivial. A single tree in the prognostic
forest that splits once on the group indicator and once on \texttt{time}
isolates the treated post-treatment cells exactly. No approximation, no
smoothing, and no large sample is required.

\subsection{This is regularisation-induced confounding, with the usual guard removed}

The phenomenon is the one Bayesian Causal Forests was designed to prevent.
\citet{Hahn2020} construct BCF precisely so that the treatment indicator does
\emph{not} enter the prognostic function; instead $\mu$ receives an estimate of
the propensity score, $\mu(\Xb, \hat\pi(\Xb))$, which supplies the information
about selection that $\mu$ needs without handing it the ability to reproduce
the treatment assignment itself. The submitted paper reverses both halves of
that prescription: the group indicator is placed inside $\mu$, and the
propensity is omitted, with the justification (Section 3.3 of the paper) that
$\hat\pi$ has been found unimportant for BCF's predictive performance in
ablation studies. Predictive unimportance and identification are different
questions. In the original BCF the propensity is not there to predict; it is
there to be the only channel through which treatment-related information
reaches $\mu$, so that the treatment indicator can be excluded.

\subsection{Parallel trends is the missing restriction}

There is a more useful way to see what has gone wrong, and it points directly
at the fix.

The paper's own description of $\mu$ is that it "captures the expected outcome
trajectory for units in the absence of treatment $\ldots$ while allowing for
\emph{static} differences between the eventually-treated and control groups".
That word is the whole of the parallel-trends assumption. Conditional parallel
trends says that, given covariates, the untreated potential outcome of the
treated group and that of the control group move together; the two groups may
sit at different levels, but the group may not interact with time. Formally,
the assumption is that $\mu$ is additively separable in the sense
\begin{equation}\label{eq:pta}
  \mu(D_i, t, \Xb_{it}) = a(D_i, \Xb_{it}) + b(t, \Xb_{it}),
\end{equation}
so that the group enters only through a level and the time path is common
across groups conditional on $\Xb$.

Equation~\eqref{eq:model} as implemented does not impose \eqref{eq:pta}. A
single unrestricted BART over $(D_i, t, \Xb)$ can represent arbitrary
group-by-time interaction in the untreated outcome, and a group-by-time
interaction in $\mu$ is observationally indistinguishable from a treatment
effect. The model therefore does not encode the identifying assumption that the
DiD framework rests on, and the ATT is not identified within it. Under
\eqref{eq:pta} the flat direction \eqref{eq:flat} disappears: $a(D_i,\Xb)$ has
no time argument and $b(t,\Xb)$ has no group argument, so their sum cannot
represent $\tau(\Xb,k)\mathbb{I}(D_i=1)\mathbb{I}(t\geq g)$, and $c$ is pinned
at zero.

This also explains why the symptom is a bias that \emph{grows} with the sample
size rather than one that shrinks. Nothing about the flat direction is a
finite-sample matter. What changes with $N$ is only how sharply the prognostic
forest can resolve the group-by-time cells against its own shrinkage prior:
with more units per cell the leaf posteriors concentrate, $\mu$ takes a larger
share, and the estimator drifts further from the truth. An estimator whose bias
increases in $N$ is not noisy, it is inconsistent.

% =============================================================================
\section{Is the problem present in the original paper and code?}
% =============================================================================

The specification is, in both, and by design rather than by accident. The
\emph{symptom} is not: the published results are unaffected, for reasons that
turn out to be informative.

\subsection{In the submitted manuscript}

Equation~(1) of \texttt{JBES\_Submission/main.tex} (page~6) defines the
prognostic function as $\mu(D_i, t, \Xb_{it})$, with the group indicator and
calendar time as explicit arguments. Equation~(4) (page~9), the estimated
specification, repeats them. The sentence immediately following Equation~(4)
states that
$\hat\pi(\Xb_i, G_i)$ is deliberately not used. The manuscript therefore
specifies exactly the configuration analysed in Section~1: the group indicator
inside $\mu$, calendar time inside $\mu$, and no propensity. Nothing in the
manuscript restricts $\mu$ to the additively separable form \eqref{eq:pta}
that its own prose describes.

\subsection{In the original simulation code}

Every notebook under \texttt{Simulation\_Studies/Original/DiD\_BCF/} uses the
same design matrix and the same forest restrictions. Taking
\texttt{DiD\_BCF\_ATE\_lin\_1.ipynb} as representative,

\begin{lstlisting}[language=Python]
data_linear['pi_hat'] = 0.5
data_linear['D'] = data_linear['post_treatment'] * data_linear['treated_group']
X = np.array(data_linear[["treated_group"] + x_columns + ["time"]])
prognostic_forest_params = {"keep_vars": np.array([0, 1] + list(range(2, 8)) + [8])}
treatment_effect_forest_params = {"keep_vars": np.array([8])}
bcf_model.sample(X_train=X, Z_train=Z, y_train=y,
                 pi_train=np.array(data_linear['pi_hat']), ...)
\end{lstlisting}

Column~0 is \texttt{treated\_group} and column~8 is \texttt{time}, and the
prognostic \texttt{keep\_vars} retains both. The propensity passed to the
sampler is the constant $0.5$, which carries no information, so the guard that
BCF would normally provide is switched off. The GATE and CATE notebooks differ
only in which columns the \emph{treatment} forest may split on; in all of them
the prognostic forest keeps the ever-treated indicator together with time. The
revision code inherited the specification unchanged, in the constants
\texttt{PROGNOSTIC\_COLS} and \texttt{TREATMENT\_COLS} of
\texttt{did\_bcf\_revision/did\_bcf.py}, where the group indicator appears twice
(as \texttt{eventually\_treated} and again as \texttt{treatment\_group}).

\subsection{Does it bite in the original design? No.}

The flat direction exists in the published specification, but whether the
posterior actually travels along it is an empirical question, and the answer
for the original simulation study is that it does not.

Reproducing the published configuration exactly (the design matrix, the forest
restrictions and the constant propensity of the code listing above, on the
original \texttt{generate\_did\_data} DGP) and reading the treatment forest
directly, the share of the true effect retained by $\hat\tau$ on treated
post-treatment rows is flat in $N$ and indistinguishable from one:

\begin{table}[htbp]
\centering
\caption{Share of the true effect retained by the treatment forest,
$\hat\tau/\att$, on treated post-treatment rows. Upper panel: the published
specification on the original DGP. Lower panel: the same estimator family on
the revision's canonical DGP, which adds an unobserved unit effect
$\alpha_i$ correlated with assignment. Replication counts in parentheses.}
\label{tab:leak}
\begin{tabular}{lccc}
\toprule
$N$ & 200 & 800 & 1600 \\
\midrule
\multicolumn{4}{l}{\emph{Original design (random assignment, no $\alpha_i$, time-varying covariates)}}\\
$\hat\tau/\att$ & 1.028 \; (3) & 1.009 \; (2) & 1.008 \; (2) \\[4pt]
\multicolumn{4}{l}{\emph{Revision canonical design (selection on an unobserved unit effect)}}\\
$\hat\tau/\att$ & 0.95,\; 0.50 & 0.47,\; 0.36 & 0.49 \\
Attenuation implied by the sweep bias & 0.83 & 0.50 & 0.43 \\
\bottomrule
\end{tabular}
\end{table}

This matters for how the revision should describe the problem, and it is a
better story than the one the diagnosis first suggested. The published results
are not merely defensible at the reported panel size; the estimator shows no
leakage at all in the design it was published on, at any $N$ tested. What
activates the flat direction is the canonical design that Reviewers~1.4
and~2.3 asked for.

The mechanism is plausible although it has not been isolated, and the
difference between the two designs is not subtle. In the original DGP,
treatment is assigned at random, there is no unobserved unit effect, and the
group indicator carries only a constant level shift of $0.75$, so the
prognostic forest has almost no reason to split on it; the covariates, drawn
per unit-period, give $\mu$ ample legitimate structure to fit instead. In the
revision DGP, $\alpha_i$ has unit standard deviation and is correlated with
assignment, so the treated and control groups differ substantially in level and
$\mu$ \emph{must} split on the group indicator to fit at all. Once a tree is
already split on the group, capturing the treatment effect costs one further
split on \texttt{time}. On this reading the unobserved confounder is what
converts a latent identification failure into an active one, which is a
testable claim: setting \texttt{alpha\_sd = 0} and \texttt{conf\_strength = 0}
in the revision DGP should restore $\hat\tau/\att \approx 1$.

The lower panel is the revision's own measurement, and there the attenuation
implied by the observed bias of the corrected estimator over the sweep,
$1 + \mathrm{bias}/\tau$, reproduces the directly measured retention to within
replication noise. The correction is not introducing that bias; it is
faithfully propagating a split between $\mu$ and $\tau$ that the likelihood
never identified.

The summary for the revision is therefore that the published numbers stand, and
that the specification nevertheless needs the restriction of Section~3 before
the paper can make any claim under selection on unobserved fixed effects. That
is a considerably more comfortable position than a correction notice, and it
turns the reviewers' request into the reason the weakness was found: the harder
canonical design they asked for is exactly what exposed it.

% =============================================================================
\section{The fix}
% =============================================================================

The requirement is to make the prognostic function unable to represent a
group-by-time interaction, while leaving it able to represent a static group
level difference, covariate-dependent trends, and arbitrary non-linearity in
the covariates. Three routes are available and they are not equally good.

The first, and the one recommended here, is to impose \eqref{eq:pta}
structurally by moving the group level out of the forest. Fit the prognostic
forest on $(\Xb, t)$ only, with the group indicator removed from its split set,
and carry the static group difference as an additive random-effect intercept.
In \texttt{stochtree} this is a single argument:

\begin{lstlisting}[language=Python]
PROGNOSTIC_COLS = ["X1", "X2", "X3", "X4", "X5", "time"]   # no group indicator
model.sample(
    X_train=X, Z_train=Z, y_train=y,
    rfx_group_ids_train=df["treatment_group"].to_numpy(dtype=np.int32),
    prognostic_forest_params={"keep_vars": prog_idx},
    treatment_effect_forest_params={"keep_vars": treat_idx},
)
\end{lstlisting}

The fitted prognostic surface is then $f(\Xb, t) + a_{G_i}$, which is exactly
\eqref{eq:pta}: the forest sees time but not the group, the random effect sees
the group but not time, and no combination of the two can produce
$\mathbb{I}(D_i{=}1)\mathbb{I}(t \geq g)$. Covariate-dependent trends
$s(\Xb)\,t$ remain expressible inside the forest, so conditional parallel
trends is imposed rather than assumed away. The cost is that group-specific
\emph{non-linear} level effects must now come through $\Xb$ rather than through
free interaction with the group indicator, which is precisely the restriction
the DiD framework asks for.

The second route is the \citet{Hahn2020} prescription: drop the group indicator
from $\mu$ and add the estimated propensity $\hat\pi(\Xb_i, G_i)$ in its place.
This is a real improvement over the status quo and it is what BCF was built to
do, but it is a mitigation rather than a guarantee. When selection is sharp,
$\hat\pi$ separates the groups well, and a tree that splits on
$\hat\pi > \text{threshold}$ and again on \texttt{time} can approximate the
interaction. It protects well in the settings where selection runs through
unobservables and less well in \texttt{B1\_selection\_obs}. It composes with
the first route and both can be used together.

The third route is to sidestep the levels entirely by fitting the model in
first differences, so that the unit effect and the static group difference are
removed before estimation. This is closest to what the posterior correction
already does internally and it would make the correction's nuisance estimate
$m^s$ coherent with the model that produced it, but it is a larger change to
the estimator than the revision needs.

Whichever route is taken, the acceptance test is the same and it should become
a permanent part of the suite: on \texttt{B2\_sweep}, the bias of both plain and
corrected DiD-BCF must \emph{decrease} in $N$. That single check would have
caught this before publication.

% =============================================================================
\section{Two repairs to the simulation design}
% =============================================================================

\subsection{The linearity-degree axis}

The axis was vacuous as delivered. In the revision DGPs the covariates are
drawn once per unit, deliberately, so that $\alpha_i$ and $f(\Xb_i)$ are
genuinely time-invariant and the design tests selection on unobserved
\emph{fixed} heterogeneity. The linearity degree entered only through
$f(\Xb_i)$, a unit constant, and through $\gamma_t$, a time constant, and both
are removed by any two-way within transformation. On a fixed seed the demeaned
outcome was identical across degrees to $2\times10^{-15}$, and
Callaway--Sant'Anna, Gardner and synthetic DiD returned bit-identical estimates
at $d=1,2,3$.

It is worth recording that this was a regression relative to the original
study, and that the cause was a deliberate and correct design choice. In
\texttt{Simulation\_Studies/Original/}, covariates are drawn per unit-period,
so the non-linear covariate contribution is time-varying and survives
differencing: on a fixed seed the two-way demeaned outcome moves by up to
$47.8$ across degrees and the naive $2\times2$ DiD moves from $3.35$ to $2.93$
to $3.97$. Drawing covariates once per unit is what the canonical design
requires; it simply removed the only channel the degree had.

The repaired axis puts the non-linearity where a differencing estimator must
confront it. The data-generating process is now written entirely in latent
covariates and does not vary with the degree at all. What varies is the
parameterisation handed to the estimators.

At $d=1$ the estimators observe the latent covariates, so the trend, the
propensity and the prognostic level are all linear in the data they see and
linear adjustment is correctly specified. At $d=2$ the three nuisance
covariates are replaced by their \citet{KangSchafer2007} transforms, as used
for the misspecified designs of \citet{SantAnna2020}: a cubic of an
interaction, $\exp(X_4/2)$, and the square of a sum. The outcome and every true
estimand are bit-identical to $d=1$; only the difficulty of recovering the
nuisance functions changes, so any difference in estimator performance is
misspecification and nothing else. At $d=3$ the effect surface additionally
becomes quadratic in $X_2$ with an $X_1$-by-$X_2$ interaction, so $\tau$ is
non-linear and non-additive in the modifiers, with $\mathrm{Var}(\tau)$ held
fixed so that the CATT-surface metrics stay comparable across the axis. The
effect modifiers $X_1$ and $X_2$ are never transformed, which keeps
"non-linear nuisance" and "non-linear effect" as separate axes.

The axis now has measurable content. Regressing the untreated outcome change on
the observed covariates and a treatment dummy, the residual violation of
conditional parallel trends that linear adjustment cannot remove is $-0.001$ at
$d=1$ and $+0.037$ at $d \geq 2$; the treatment-effect surface is identical at
$d=1$ and $d=2$ and differs by up to $10.5$ at $d=3$, with
$\mathrm{SD}(\tau)$ moving only from $1.063$ to $1.069$. These are asserted by
\texttt{Results/\_staging/check\_linearity\_axis.py}, which should be run
whenever the DGP module is touched.

One consequence should be stated plainly. Under
\texttt{selection="unobservable"}, which is the setting of
\texttt{B1\_baseline}, \texttt{B1\_strong\_confounder},
\texttt{B1\_serial\_corr}, \texttt{B1\_null} and both sweep scenarios, the
covariates are balanced across arms by construction, so no adjustment at all is
needed and the linearity degree loads only through nuisance-fitting variance.
It moves \emph{bias} only where assignment depends on the covariates. That was
true of \texttt{B1\_selection\_obs} alone, so a scenario
\texttt{B1\_selection\_both} has been added, combining an unobserved confounder
with covariate-driven assignment. It is the only cell of the design in which
conditional parallel trends is genuinely violated, and it is where the repaired
axis does its work.

\subsection{The pre-trend recentring}

\texttt{plain\_estimands} defaulted to reporting
$\hat\tau(i,t) - \hat\tau(i,g-1)$, on the documented reasoning that under
conditional parallel trends the subtracted term is approximately zero. The
premise is right and the conclusion does not follow. Because $Z = D_{it}$ is
zero on every pre-treatment row, those rows contribute nothing to the
likelihood of $\tau$; the treatment forest has no gradient with which to split
on time before adoption, and $\hat\tau(i, g-1)$ is not an estimate of a
pre-treatment effect but the post-treatment estimate re-evaluated at a
different time value. The subtraction removes the entire estimate rather than a
small correction. Measured on one \texttt{B1\_baseline} replication,
$\hat\tau$ is $3.655$ on treated post rows and $3.669$ at $g-1$ against a true
ATT of $3.849$, so the reported ATT is $3.655$ with the flag off and $-0.014$
with it on; the same near-equality holds at every panel size tested.

The default is now \texttt{False} and the flag is retained only to reproduce
the earlier runs. The command-line switch in \texttt{scripts/run\_did\_bcf.py}
has been inverted accordingly, from \texttt{--no-pretrend-recenter} to
\texttt{--pretrend-recenter}.

This leaves the revision without a pre-trend diagnostic, which
Reviewer~3.2.3 asked for, and the reason is structural rather than incidental.
A pre-trend test requires $\tau$ to be estimable before adoption, which means
the unconstrained specification with $Z = \mathbb{I}(G_i \neq \infty)$ active in
all periods. That is Equation~(1) of the manuscript, the specification the
paper argues against on complexity grounds, and the earlier experience that it
estimates treatment effects less well is consistent with that argument. The two
goals are in tension and the sensible resolution is to run the unconstrained
specification as a \emph{diagnostic} fit alongside the constrained one used for
estimation, reporting $\hat\tau(\Xb, k)$ for $k<0$ from the former and the ATT
from the latter. This costs one extra fit per replication and asks the
unconstrained model only for the thing it is good at.

% =============================================================================
\section{Status}
% =============================================================================

The code changes described in Section~4 are in place in
\texttt{did\_bcf\_revision/} and pass the smoke and acceptance tests. They do
not require re-running the R benchmarks, whose estimates depend only on the
observed panel, but they do change the panels at $d \geq 2$, so every estimator
must be re-run to produce a coherent comparison. The identification fix of
Section~3 is not yet implemented, because it changes the estimator rather than
the design and should be validated on its own before the full grid is
committed. The recommended sequence is to implement the additive
group-effect specification, run \texttt{B2\_sweep} alone at all four panel
sizes, confirm that the bias of both plain and corrected DiD-BCF decreases in
$N$, and only then re-run the remaining scenarios.

\begin{thebibliography}{9}
\bibitem[Hahn et al.(2020)]{Hahn2020}
Hahn, P. R., Murray, J. S., and Carvalho, C. M. (2020).
Bayesian regression tree models for causal inference: regularization,
confounding, and heterogeneous effects.
\emph{Bayesian Analysis} 15(3), 965--1056.
\url{https://doi.org/10.1214/19-BA1195}

\bibitem[Kang and Schafer(2007)]{KangSchafer2007}
Kang, J. D. Y. and Schafer, J. L. (2007).
Demystifying double robustness: a comparison of alternative strategies for
estimating a population mean from incomplete data.
\emph{Statistical Science} 22(4), 523--539.
\url{https://doi.org/10.1214/07-STS227}

\bibitem[Sant'Anna and Zhao(2020)]{SantAnna2020}
Sant'Anna, P. H. C. and Zhao, J. (2020).
Doubly robust difference-in-differences estimators.
\emph{Journal of Econometrics} 219(1), 101--122.
\url{https://doi.org/10.1016/j.jeconom.2020.06.003}
\end{thebibliography}

\end{document}
